\documentclass[11pt]{article}
\usepackage[T1]{fontenc}
\usepackage{lmodern}
\usepackage{amsmath,amssymb,amsthm}
\usepackage[margin=1in]{geometry}
\usepackage[colorlinks=true,linkcolor=blue,citecolor=blue,urlcolor=blue]{hyperref}
\hypersetup{pdftitle={Uniform covariance bounds for constrained Ising replicas},pdfauthor={Kaiwen Jin}}
\newtheorem{theorem}{Theorem}[section]
\newtheorem{proposition}[theorem]{Proposition}
\newtheorem{lemma}[theorem]{Lemma}
\title{Uniform covariance bounds for constrained Ising replicas}
\author{Kaiwen Jin\\\small\texttt{yui.ui.siki@gmail.com}}
\date{}
\begin{document}
\maketitle
\begin{abstract}
We prove that the covariance of finitely many Ising replicas conditioned on
fixed magnetizations has a bound independent of the number of replicas,
the conserved totals and all replica-dependent external fields, for every
fixed interaction matrix. The constraints may conserve the spin totals in
each part of an arbitrary partition. This proves the covariance conjecture
of Caputo and Morellini. The argument conditions two independent copies on
their coordinatewise sum, which cancels the external fields. The remaining
disagreements have finitely many symmetric block laws. Pairing identical
laws gives a lower bound for the conditioning probability, and a Fourier
estimate gives a second-moment numerator of the same order.
\end{abstract}

\section{The covariance question}
Magnetization constraints couple otherwise independent replicas of a spin
system. A covariance bound uniform in the replica count controls the variance
of every linear spin statistic, including in the presence of arbitrary
external fields. Such bounds arise in the study of spin-exchange dynamics
with conserved magnetizations; see Caputo and Morellini~\cite[Sections 1.5
and 4.5]{CM}.

Fix an integer $n\geq1$, a real symmetric matrix $J\in\mathbb R^{n\times n}$,
and a partition $\mathcal A$ of $[n]=\{1,\ldots,n\}$. A configuration of
$N\geq2$ replicas is
$\sigma=(\sigma(1),\ldots,\sigma(N))\in(\{-1,1\}^n)^N$.
For integers $M_A$, $A\in\mathcal A$, suppose that
\[
 \Omega_{N,M}=\left\{\sigma:
  \sum_{i=1}^N\sum_{j\in A}\sigma_j(i)=M_A
  \text{ for every } A\in\mathcal A\right\}
\]
is nonempty. For any $v=(v_{ij})\in\mathbb R^{Nn}$, define
\begin{equation}\label{eq:law}
 \nu_{N,M,v}(\sigma)=\frac{\mathbf1_{\Omega_{N,M}}(\sigma)}{Z_{N,M,v}}
 \exp\left\{\frac12\sum_{i=1}^N\sigma(i)^TJ\sigma(i)
                  +\sum_{i,j}v_{ij}\sigma_j(i)\right\}.
\end{equation}
The normalizer $Z_{N,M,v}$ is a positive finite sum.
This is the canonical replica measure of
\cite[Definition~1.9, eqs.\ (1.23)--(1.24)]{CM}, with an arbitrary linear
tilt and the conserved densities written as integer spin totals: $\rho(A)=(1+M_A/(N|A|))/2$.

Caputo and Morellini ask in \cite[Section 4.5]{CM} whether its covariance
admits a bound uniform in $N$, the feasible totals and every tilt $v$, for
arbitrary fixed $J$. Their Proposition 3.6 gives the bound
$2/(1-2\lambda_{\max}(J))$ for a single total constraint when $J$ is positive
definite and $\lambda_{\max}(J)<1/2$; Section 3.4 treats the corresponding
partitioned constraints at high temperature. We prove the uniform covariance
bound without a spectral-smallness assumption.

\begin{theorem}\label{thm:main}
For every $n\geq1$, every real symmetric $J\in\mathbb R^{n\times n}$ and every
partition $\mathcal A$ of $[n]$, there is a finite constant
$C(J,\mathcal A)$ such that
\begin{equation}\label{eq:main}
 \bigl\|\operatorname{Cov}_{\nu_{N,M,v}}(\sigma)\bigr\|_{\mathrm{op}}
 \leq C(J,\mathcal A)
\end{equation}
for every $N\geq2$, every feasible integer vector $M$, and every
$v\in\mathbb R^{Nn}$. The covariance is that of all $Nn$ spins, and the
operator norm is induced by the ordinary Euclidean norm. More explicitly,
for every real array $a$,
\begin{equation}\label{eq:variance}
 \operatorname{Var}_{\nu_{N,M,v}}\left(\sum_{i,j}a_{ij}\sigma_j(i)\right)
 \leq C(J,\mathcal A)\sum_{i,j}a_{ij}^2.
\end{equation}
\end{theorem}
For fixed $n$, there are finitely many partitions, so taking their maximum
also gives a bound depending only on $J$. The constant may grow with the fixed
dimension and interaction; no optimal dependence is sought here. The result proves the covariance conjecture in
\cite[Section 4.5]{CM}. We do not derive entropy-decay estimates in this note.

The external fields can vary freely from one replica to another, so a
conditioning estimate for one fixed tilted law does not provide the required
uniformity. Instead, take two independent configurations and condition on
their coordinatewise sum. Their difference then has no external field and
zero conserved totals. Only the subset of differing sites in each replica
matters, giving finitely many possible block laws. Pairing blocks with
identical laws makes a Fourier integral nonnegative, which yields a uniform
lower bound for the zero-total event. A centered second-moment expansion
has the same dependence on the active counts, and the two bounds cancel
upon conditioning.

Section~\ref{sec:blocks} proves the required estimate for conditioned block
sums. Section~\ref{sec:disagreement} applies it to the two-copy difference
and completes the proof. Empty subsets and parts with no disagreements are
included throughout.
\section{Conditioned block sums}\label{sec:blocks}
For a subset $D\subset[n]$, let $J_D$ be the principal submatrix of $J$ and
put
\begin{equation}\label{eq:masklaw}
 \mu_D(s)=\frac{\exp(s^T J_D s)}
 {\sum_{u\in\{-1,1\}^D}\exp(u^T J_D u)},
 \qquad s\in\{-1,1\}^D.
\end{equation}
The empty cube has one element, so $\mu_\varnothing$ is a point mass.
We extend its vectors by zero outside $D$. These measures are symmetric under
$s\mapsto-s$. Set
\begin{equation}\label{eq:delta}
 d=|\mathcal A|,\qquad B=\sum_{j<k}|J_{jk}|,\qquad \delta=e^{-4B}.
\end{equation}
The diagonal part of $s^T J_Ds$ is constant. Its remaining part has
oscillation at most $4B$, and hence
\begin{equation}\label{eq:density}
 \mu_D(s)\geq\delta\,2^{-|D|}\geq p_*:=\delta\,2^{-n}>0.
\end{equation}

Let $X_1,\ldots,X_N$ be independent with laws $\mu_{D_1},\ldots,\mu_{D_N}$,
where the subsets $D_i$ are arbitrary. We condition their sum on the same
conserved quantities as in the original problem. Write
\[
 m_{iA}=|D_i\cap A|,\qquad K_A=\sum_i m_{iA},\qquad
 S_{i,A}=\sum_{j\in D_i\cap A}X_{ij},\qquad S=\sum_i S_i.
\]
The event $S=0$ is possible exactly when each $K_A$ is even. Necessity follows
from parity. For sufficiency, assign half the active coordinates in each part
sign $+1$ and half sign $-1$; every such assignment has positive probability
by \eqref{eq:density}.

\begin{proposition}\label{prop:moment}
There is a finite $\kappa(J,\mathcal A)$ such that, for every $N\geq2$, every
sequence of subsets $D_i$ for which all $K_A$ are even, and every real array
$a=(a_{ij})$,
\begin{equation}\label{eq:moment}
 \mathbb E\left[\left(\sum_i\sum_{j\in D_i}a_{ij}X_{ij}\right)^2
       \,\middle|\, S=0\right]
 \leq \kappa(J,\mathcal A)\sum_i\sum_{j\in D_i}a_{ij}^2.
\end{equation}
\end{proposition}
The conditioning probability can tend to zero with $N$. We will bound it
from below and the corresponding second-moment numerator from above at the
same scale, including parts with few active coordinates.

For $t=(t_A)_{A\in\mathcal A}\in\mathbb R^d$, put
$\phi_i(t)=\mathbb E e^{\mathrm i t\cdot S_i}$. Symmetry makes $\phi_i$ real.
The uniform cube has characteristic function
$\prod_A(\cos t_A)^{m_{iA}}$. Subtracting $\delta$ times the uniform law
from $\mu_{D_i}$ leaves a nonnegative measure of mass $1-\delta$, so
\[
 |\phi_i(t)|\leq 1-\delta+
             \delta\prod_A|\cos t_A|^{m_{iA}}.
\]
For $c_1,\ldots,c_m\in[0,1]$ with $m\leq n$, the product is at most each
$c_\ell$, giving
$\sum_\ell(1-c_\ell)\leq n(1-\prod_\ell c_\ell)$.
Together with $1-|\cos t|\geq\frac12\sin^2t$ and $1-x\leq e^{-x}$,
this proves
\begin{equation}\label{eq:damping}
 |\phi_i(t)|\leq
 \exp\left(-h\sum_A m_{iA}\sin^2t_A\right),
 \qquad h=\frac{\delta}{2n}.
\end{equation}
The formula also holds for an empty subset $D_i$.

\begin{lemma}\label{lem:denominator}
If all $K_A$ are even, then
\begin{equation}\label{eq:denominator}
 \mathbb P(S=0)\geq L\prod_A(1+K_A)^{-1/2},
 \qquad
 L=p_*^{\,2^n}(r/\pi)^d e^{-ndr^2},\quad r=\frac1{2n}.
\end{equation}
\end{lemma}
\begin{proof}
Pair blocks with identical subsets $D_i$, leaving at most one block of each type
unpaired. There are at most $2^n$ unpaired blocks. Each pair contributes an
even active count to each part. Thus the unpaired counts are even as well,
and we can choose one unpaired spin assignment whose totals are all zero.
Its probability is at least $p_*^{2^n}$, by \eqref{eq:density}.

Let $S^{\mathrm p}$ be the sum of the paired blocks and $K_A^{\mathrm p}$
their active counts. If $2r_D$ blocks of type $D$ are retained, finite Fourier
inversion gives
\begin{equation}\label{eq:paired}
 \mathbb P(S^{\mathrm p}=0)=
 (2\pi)^{-d}\int_{[-\pi,\pi]^d}
       \prod_{D\subset[n]}\phi_D(t)^{2r_D}\,dt.
\end{equation}
Here $\phi_D$ is the common characteristic function of blocks of type $D$.
The integrand is nonnegative everywhere, because these functions are real
and their exponents are even. We may therefore retain only the box
\[
 |t_A|\leq\frac{r}{\sqrt{1+K_A^{\mathrm p}}},\qquad A\in\mathcal A.
\]
For each paired block, the pointwise Cauchy--Schwarz inequality gives
\[
 \mathbb E(t\cdot S_i)^2
 \leq n\sum_A m_{iA}t_A^2\leq n^2r^2=\frac14.
\]
Using $\cos u\geq1-u^2/2$ and
$\log(1-z)\geq-2z$ for $0\leq z\leq1/2$, we obtain on this box
\[
 \phi_i(t)\geq
 1-\frac n2\sum_A m_{iA}t_A^2
 \geq\exp\left(-n\sum_A m_{iA}t_A^2\right).
\]
The product in \eqref{eq:paired} is consequently at least
$e^{-n\sum_A K_A^{\mathrm p}t_A^2}\geq e^{-ndr^2}$ there. Its integral
is bounded below by
\[
 (r/\pi)^d e^{-ndr^2}\prod_A(1+K_A^{\mathrm p})^{-1/2}.
\]
The chosen unpaired assignment is independent of the paired blocks.
Multiplying its probability by this bound, then using
$K_A^{\mathrm p}\leq K_A$, proves \eqref{eq:denominator}.
Empty products are one, so the argument includes the case with no pairs.
\end{proof}

\begin{lemma}\label{lem:numerator}
For all even counts $K_A$ and all real arrays $a$,
\begin{equation}\label{eq:numerator}
 \mathbb E\left[F^2\mathbf1_{\{S=0\}}\right]
 \leq U E_a\prod_A(1+K_A)^{-1/2},
\end{equation}
where
\[
 F=\sum_i\sum_{j\in D_i}a_{ij}X_{ij},\qquad
 E_a=\sum_i\sum_{j\in D_i}a_{ij}^2,
\]
and the following constants depend only on $J$ and $\mathcal A$:
\begin{equation}\label{eq:upperconstants}
 \gamma=\frac{2\delta}{n\pi^2},\qquad
 H=4\sqrt{\frac{2n}{\delta}},\qquad
 U=e^\delta\left(n+\frac{2n^2}{\gamma}\right)H^d.
\end{equation}
\end{lemma}
\begin{proof}
Put $F_i=\sum_{j\in D_i}a_{ij}X_{ij}$,
$E_i=\sum_{j\in D_i}a_{ij}^2$, and $Q(t)=\sum_A K_A t_A^2$.
Symmetry gives $\mathbb EF_i=0$. Cauchy--Schwarz gives
\[
 \mathbb EF_i^2\leq nE_i,\qquad
 \mathbb E(t\cdot S_i)^2\leq n\sum_A m_{iA}t_A^2.
\]
Define $B_i(t)=\mathbb E[F_i e^{\mathrm i t\cdot S_i}]$.
The vanishing mean and $|e^{\mathrm iu}-1|\leq|u|$ yield
\[
 |B_i(t)|
 =\left|\mathbb E[F_i(e^{\mathrm i t\cdot S_i}-1)]\right|
 \leq n\sqrt{E_i}\left(\sum_A m_{iA}t_A^2\right)^{1/2}.
\]
Summing and applying Cauchy--Schwarz over blocks, we get
\begin{equation}\label{eq:firstmoments}
 \left(\sum_i|B_i(t)|\right)^2\leq n^2 E_a Q(t).
\end{equation}

Let $V_i(t)=\mathbb E[F_i^2 e^{\mathrm i t\cdot S_i}]$.
Independence gives the exact expansion
\begin{align}\label{eq:expansion}
 \mathbb E[F^2e^{\mathrm i t\cdot S}]
 &=\sum_i V_i(t)\prod_{k\ne i}\phi_k(t)\notag\\
 &\quad+\sum_{i\ne j}B_i(t)B_j(t)
                         \prod_{k\notin\{i,j\}}\phi_k(t).
\end{align}
Each term omits at most two characteristic factors. For any omitted index
set $I$ with $|I|\leq2$, \eqref{eq:damping} implies
\begin{align*}
 \prod_{k\notin I}|\phi_k(t)|
 &\leq\exp\left(-h\sum_AK_A\sin^2t_A
       +h\sum_{i\in I,A}m_{iA}\sin^2t_A\right)\\
 &\leq e^{2nh}\exp\left(-h\sum_AK_A\sin^2t_A\right).
\end{align*}
The omitted blocks contain at most $2n$ active coordinates. On
$[-\pi/2,\pi/2]^d$ we have $\sin^2t_A\geq4t_A^2/\pi^2$ and $2nh=\delta$,
so the last expression is at most $e^\delta e^{-\gamma Q(t)}$.
Since $|V_i(t)|\leq nE_i$ and
$\sum_{i\ne j}|B_iB_j|\leq(\sum_i|B_i|)^2$,
\eqref{eq:expansion} and \eqref{eq:firstmoments} give
\begin{equation}\label{eq:integrand}
 \left|\mathbb E[F^2e^{\mathrm i t\cdot S}]\right|
 \leq e^\delta E_a[n+n^2Q(t)]e^{-\gamma Q(t)}.
\end{equation}

Every $S_A$ is an even integer. Hence
\begin{equation}\label{eq:smalltorus}
 \mathbb E[F^2\mathbf1_{\{S=0\}}]
 =\pi^{-d}\int_{[-\pi/2,\pi/2]^d}
                    \mathbb E[F^2e^{\mathrm i t\cdot S}]\,dt.
\end{equation}
Both this formula and \eqref{eq:paired} follow by integrating each
exponential monomial: the normalized integral is one for frequency zero
and zero for any other allowed integer frequency. Only finite sums occur.

For $q\geq0$,
\[
 [n+n^2q]e^{-\gamma q}
 \leq\left(n+\frac{2n^2}{\gamma}\right)e^{-\gamma q/2}.
\]
Write $\alpha=\gamma/2=\delta/(n\pi^2)$ and recall
$H=4\sqrt{2n/\delta}$. For an integer $k\geq1$, the substitution $u=\sqrt{\alpha k}\,t$ gives
\[
 \frac1\pi\int_{-\pi/2}^{\pi/2}e^{-\alpha k t^2}\,dt
 \leq\frac4{\pi\sqrt{\alpha k}}
 \leq\frac{H}{\sqrt{1+k}}.
\]
Indeed, the integral of $e^{-u^2}$ over any finite symmetric interval is
at most $2+2e^{-1}<4$: use $e^{-u^2}\leq e^{-|u|}$ when $|u|\geq1$.
For $k=0$ the normalized integral is one, which is also bounded by $H$.
Applying these one-dimensional bounds in \eqref{eq:smalltorus}, with
\eqref{eq:integrand}, proves \eqref{eq:numerator}.
\end{proof}

\begin{proof}[Proof of Proposition~\ref{prop:moment}]
Divide \eqref{eq:numerator} by the positive lower bound
\eqref{eq:denominator}. The factors
$\prod_A(1+K_A)^{-1/2}$ cancel, and \eqref{eq:moment} follows with
$\kappa(J,\mathcal A)=U/L$.
\end{proof}

\section{Disagreements of two independent copies}\label{sec:disagreement}
\begin{proof}[Proof of Theorem~\ref{thm:main}]
Let $\sigma,\tau$ be independent with law $\nu_{N,M,v}$, and set
\[
 c_{ij}=\frac{\sigma_j(i)+\tau_j(i)}2,\qquad
 x_{ij}=\frac{\sigma_j(i)-\tau_j(i)}2.
\]
Fix a value of $c$ with positive probability. Where $c_{ij}\ne0$, the
coordinate $x_{ij}$ is zero; where $c_{ij}=0$, it is $+1$ or $-1$.
Put $D_i=\{j:c_{ij}=0\}$. Because both copies have total $M_A$,
\[
 \sum_i\sum_{j\in A}c_{ij}=M_A,\qquad
 \sum_i\sum_{j\in A}x_{ij}=0 \quad(A\in\mathcal A).
\]
Conversely, for this fixed $c$, the zero totals of $x$ are precisely the
conditions that $c+x$ and $c-x$ both belong to $\Omega_{N,M}$.

The pair energy in a block is
\begin{align*}
 &\tfrac12(c_i+x_i)^T J(c_i+x_i)
   +\tfrac12(c_i-x_i)^T J(c_i-x_i)\\
 &\hspace{2em}+v_i^T(c_i+x_i)+v_i^T(c_i-x_i)
   =c_i^T Jc_i+x_i^T Jx_i+2v_i^Tc_i.
\end{align*}
The terms involving $v$ and $c$ are constant under this conditioning.
The map from an allowed orientation $x$ to the ordered pair
$(c+x,c-x)$ is one-to-one. Thus the conditional law of $x$ is exactly the
product of the laws \eqref{eq:masklaw} for the subsets $D_i$, conditioned on
$S=0$. In particular, all its active counts $K_A$ are even.

For $F_a(\sigma)=\sum_{i,j}a_{ij}\sigma_j(i)$, the independent-copy identity
and Proposition~\ref{prop:moment} now give
\begin{align*}
 \operatorname{Var}_{\nu_{N,M,v}}F_a
 &=\tfrac12\mathbb E\bigl(F_a(\sigma)-F_a(\tau)\bigr)^2
   =2\mathbb E\left(\sum_{i,j}a_{ij}x_{ij}\right)^2\\
 &\leq\frac{2U}{L}\,
       \mathbb E\sum_i\sum_{j\in D_i}a_{ij}^2
 \leq\frac{2U}{L}\sum_{i,j}a_{ij}^2.
\end{align*}
All conditional laws were taken only at supported values of $c$.
The constants $L>0$ and $U<\infty$ depend on $n,J,\mathcal A$ alone.

To obtain the operator bound, let $\Sigma=\operatorname{Cov}_{\nu_{N,M,v}}(\sigma)$.
For arbitrary real $a,b$, Cauchy--Schwarz gives
\[
 |b^T\Sigma a|\leq\sqrt{\operatorname{Var}F_b\operatorname{Var}F_a}
 \leq(2U/L)\|a\|_2\|b\|_2.
\]
For $\Sigma a\ne0$, take $b=\Sigma a/\|\Sigma a\|_2$; the case
$\Sigma a=0$ is immediate. This proves \eqref{eq:main} with
$C(J,\mathcal A)=2U/L$.
\end{proof}
In terms of $n$, $d=|\mathcal A|$ and
$\delta=\exp(-4\sum_{j<k}|J_{jk}|)$, one explicit choice is
\begin{align*}
 C(J,\mathcal A)
 &=2e^{\delta+d/(4n)}
 \left(n+\frac{n^3\pi^2}{\delta}\right)
 \left(\frac{2^n}{\delta}\right)^{2^n}\\
 &\quad\times\left(8n\pi\sqrt{\frac{2n}{\delta}}\right)^d.
\end{align*}
It is finite for every fixed real symmetric $J$. Its large size reflects the
rough estimates above, rather than an asserted sharp dependence on the
interaction or dimension.
\paragraph{AI assistance.}
OpenAI Codex assisted with mathematical exploration, proof development and
checking, literature search, and LaTeX drafting. Independent mathematical
review is requested.

\begin{thebibliography}{1}
\bibitem{CM}
P. Caputo and M. Morellini,
\emph{Kac's Program and Relative Entropy Decay for Nonlinear Spin-Exchange
Dynamics}, arXiv:2511.05223v1 (2025).
\href{https://arxiv.org/abs/2511.05223v1}{arXiv:2511.05223v1}.
\end{thebibliography}
\end{document}
